% Standalone section; ordinary mathematical proof and a rejected compilation.
\section{An inductive seed invariant and a specialized halt test}
\label{sec:seed-invariant}

The exact blank-start calibration supplies logical registers
\[
 (R_0,\ldots,R_8)=(3,0,0,0,3,0,0,0,0).
\]
This section records a sound way to exploit that seed.  It also records why
one promising source optimization has not yet produced an accepted binary
machine.  The rejected table must not be counted as a construction result.

\subsection{Finite abstraction with an unbounded meaning}

Fix a cutoff $C\geq4$ and modulus $M\geq1$.  Partition the naturals into
the exact values $0,\ldots,C-1$ and the $M$ sets
\[
 B_j=\{n\geq C:n\equiv j\pmod M\}\qquad(0\leq j<M).
\]
A mask is a union of these buckets.  At each instruction occurrence the
certificate gives a Cartesian product of nine masks, or the empty set.
Its meaning is an overapproximation of all reachable register valuations.
No upper bound on a register, execution time, or number of source iterations
is assumed.

Increment and positive decrement act on masks by their exact bucket images.
These images can be computed from the representatives in
$[C,C+2M)$, together with the exact values below $C$.  Beyond the cutoff,
increment changes only the residue.  Decrement does likewise except at
$n=C$, where it reaches the exact value $C-1$.  Including the next
representative of that same residue handles both outcomes.  Thus this
finite computation is justified for every natural number.

At a decrement instruction, the zero edge intersects the selected register
with $\{0\}$; the positive edge intersects it with the positive naturals
and then applies decrement.  Increment changes just its selected register.
Break and continue retain every register and follow their exact path
successor.  The checker reconstructs every successor by binary path
increment and suffix removal and rejects a nonzero unread suffix.

The acceptance conditions are elementary: the initial masks contain the
calibrated seed, and every abstract transition image is included in its
target masks.  Induction on a finite source execution now proves the stated
invariant.  A separate verification pass checks these conditions after the
worklist computation has finished; it can also check an existing certificate
without running that computation.  The standard-library-only program
\texttt{portable\_seed\_invariant.py} implements both modes.

For the retained DN4 source, $C=8$, $M=2$ gives $4745$ instruction
occurrences and $6162$ nonempty edge obligations.  All obligations pass.
There are three empty instruction masks and two decrement guards whose
positive edges are impossible.  In particular, at source path
\texttt{11101001010}, the unique initial increment of register $6$ in
the NIP routine, its incoming mask is exactly $R_6=0$.  Its source
index is $4599$ when indices start at zero.  The corresponding original
source index is $4975$.  These facts are certified by the exact source
hash recorded in \texttt{dn4\_seed\_invariant.json}.

\subsection{A correct source specialization}

Write $D_r$ for destructive decrement, $I_r$ for increment, and $H$ for
halt.  A binary pair is the original structured register-machine cons
cell: a zero decrement in its left child exits that cell normally,
and a break or continue in its right child is handled by the inherited
structured rules.  Let
\[
 U=(D_5,((D_5,\mathrm{continue}\ 1),\mathrm{break}\ 1)).
\]
This is the source routine hosted at normalized state $88$.  If its
input register $R_5$ is zero, it returns normally.  If $R_5>0$, it
clears $R_5$ and exits with one break.  The inner loop decreases $R_5$
on every iteration, so both claims hold for arbitrary naturals.

The original NIP halt test increments the initially zero scratch register
$R_6$ and invokes the general destructive equality test on $R_5,R_6$.
That test clears both registers and continues to halt precisely when they
were equal.  Hence the complete original fragment halts precisely when
the incoming $R_5=1$; otherwise it returns with $R_5=R_6=0$.

Replace that fragment by
\[
 P=(D_5,(U,H)).
\]
For incoming $R_5=0$, the first decrement exits $P$, leaving zero.
For $R_5=1$, it leaves zero, $U$ returns normally, and $H$ is executed.
For $R_5\geq2$, it leaves a positive value, so $U$ clears the value
and breaks past $H$.  Register $R_6$ is unchanged and therefore remains
zero.  No other register changes.  Thus $P$ has exactly the original
closed-fragment outcome, for every input satisfying $R_6=0$.

Both fragments terminate, and the three cases above prove their full
effect, not just agreement on sampled counter values.  The inductive
entry invariant allows their substitution in the surrounding source
program from the calibrated seed.  The binary interpreter may append a
different number of spare blocks because fewer source increments occur;
this does not change register-machine semantics, and the universal
runtime theorem quantifies over arbitrary nonempty finite spare lists.

\subsection{Why the direct binary patch is rejected}

The source replacement shortens the original $5121$-leaf source to
$5115$ leaves, or the DN4 source from $4745$ to $4739$ leaves.
A direct attempt reuses normalized states $38,53,68$:
state $38$ becomes the node $(P,\mathcal D_{69})$, state $53$ the node $P$,
and state $68$ the node $(U,H)$.  Five transition rows change, and the
physical state count remains $225$. Here $\mathcal D_{69}$ denotes the
unchanged source subtree hosted at normalized state 69.

The source decoder alone accepts this layout, but the complete binary
runtime rejects it.  In particular, normalized row
\[
 (38,0)\longmapsto(0,\mathrm R,13)
\]
also serves as an interpreter helper while executing the unrelated
increment-$8$ instruction at path \texttt{01010}.  The attempted
specialization changes its target to $53$.  Having a unique logical
decoder occurrence does not imply exclusive use as a physical Turing
state.

There is a direct finite counterexample.  Start the increment-$8$
primitive at the normal dispatcher, with header width $64$, all nine
primary registers zero, one zero-valued spare block, and blank exterior.
The inherited table reaches the correct complete output frame after
$636$ transitions: the path is \texttt{01011}, register $8$ is one,
and a one-valued spare block has been appended.  The attempted table
instead reaches a complete frame at that same next path after $648$
transitions, with all primary registers still zero and its original
single spare unchanged.  Thus it has performed a no-op in place of
the required increment.  The complete input and output supports are
recorded in \texttt{dn4\_nip\_rejection.json}.  The standalone checker
\texttt{verify\_nip\_rejection.py} reconstructs the five-row patch and
replays this counterexample directly from the accepted DN4 table.

This is also relevant immediately to the calibrated computation:
the increment at \texttt{01010} is the third source instruction from
the seed, following a zero test of register $8$ and a positive decrement
of register $0$.  In addition, a blank-start diagnostic finds $99$
uses of the affected helper row in the DN4 startup.  The changed
table therefore has neither the required universal runtime certificate
nor an inherited startup certificate.

The source specialization and seed invariant remain usable research
results.  A successful compilation must preserve or replace the
interpreter helper, as well as the logical decoder edges.
